\section{Ingestion Pipeline}
\label{sec:ingestion}

\subsection{Overview}
This section describes the low-rate JSON path (\texttt{microgrid/\#}), which is fully configuration-driven via Telegraf's MQTT consumer input plugin, regex and Starlark processors, and the InfluxDB v2 output plugin, with no custom code. Telegraf runs as a Docker container on the historian server alongside Mosquitto and InfluxDB. The full-rate binary stream is ingested by a separate service (Section~\ref{sec:fullrate}).

\subsection{Telegraf Configuration}
The as-built configuration (\path{telegraf/telegraf.conf}) is:
\begin{lstlisting}[language=toml]
[agent]
  interval = "10s"
  flush_interval = "1s"
  precision = "1ns"
  omit_hostname = true

[[inputs.mqtt_consumer]]
  servers = ["tcp://mosquitto:1883"]
  topics = ["microgrid/#"]
  qos = 0
  data_format = "json"
  json_time_key = "ts"
  json_time_format = "unix_ns"
  topic_tag = "topic"

[[processors.regex]]
  order = 1
  [[processors.regex.tags]]
    key = "topic"
    pattern = "^microgrid/([^/]+)/([^/]+)/([^/]+)$"
    replacement = "${1}"
    result_key = "node_id"
  # ... likewise ${2} -> subsystem, ${3} -> metric

[[outputs.influxdb_v2]]
  urls = ["http://influxdb:8086"]
  token = "${INFLUXDB_TOKEN}"
  organization = "qut-microgrid"
  bucket = "telemetry"
\end{lstlisting}
Notes:
\begin{itemize}[noitemsep]
  \item \texttt{json\_time\_key = "ts"} and \texttt{json\_time\_format = "unix\_ns"} make Telegraf use the node's PTP-derived timestamp as the InfluxDB point timestamp verbatim, and \texttt{precision = "1ns"} preserves its full resolution.
  \item The full topic is kept as the \texttt{topic} tag, and the regex processor extracts \texttt{node\_id}, \texttt{subsystem}, and \texttt{metric} from it. This replaced the \texttt{topic\_parsing} block of the original design.
  \item Telegraf reaches the broker and database by container name over the internal Docker network.
  \item The InfluxDB token is injected via an environment variable and is never stored in the configuration file. The broker requires no credentials (see the MQTT design document).
\end{itemize}

\subsection{Burst Payload Path}
Each JSON burst \texttt{\{"ts", "sync", "dt\_ns", "v": [20 values]\}} is parsed into a single point: Telegraf's JSON parser flattens the \texttt{v} array into fields \texttt{v\_0} to \texttt{v\_19}, and \texttt{sync} and \texttt{dt\_ns} become numeric fields. The parser cannot reconstruct per-sample timestamps itself. Booleans and strings are dropped by the parser, which is why the firmware sends \texttt{sync} as an integer.

\paragraph{Burst Splitting} A Starlark processor, run after the regex processor, splits each burst into individually timestamped points. For every burst point it emits 20 additional points in the \texttt{samples} measurement, sample $k$ at \texttt{ts} $+ k \cdot$ \texttt{dt\_ns}, carrying the burst's tags and the fields \texttt{v} and \texttt{sync}:

\begin{lstlisting}[language=toml]
[[processors.starlark]]
  order = 2
  source = '''
def apply(metric):
    if "dt_ns" not in metric.fields or "v_0" not in metric.fields:
        return metric
    dt = int(metric.fields["dt_ns"])
    sync = metric.fields.get("sync")
    out = [metric]
    for k in range(1024):
        key = "v_" + str(k)
        if key not in metric.fields:
            break
        sample = Metric("samples")
        for name, value in metric.tags.items():
            sample.tags[name] = value
        sample.fields["v"] = metric.fields[key]
        if sync != None:
            sample.fields["sync"] = sync
        sample.time = metric.time + k * dt
        out.append(sample)
    return out
'''
\end{lstlisting}

The original burst point is retained alongside the samples, so that existing dashboards and the liveness check are unaffected; at the JSON path's rate the duplication is negligible. This path is the low-rate fallback for when the full-rate stream is not in use. The full-rate binary stream on \texttt{daq2/<node\_id>/samples} lies outside Telegraf's subscribed topic tree and is ingested separately (Section~\ref{sec:fullrate}).

\subsection{Write Batching}
\begin{table}[H]
  \caption{Write Batching}
  \begin{center}
    \begin{tabular}{lll}
      \hline
      \textbf{Parameter} & \textbf{Value} & \textbf{Notes} \\
      \hline
      \texttt{metric\_batch\_size} & 1000 points (default) & Maximum points per write request \\
      \texttt{flush\_interval} & 1s & Maximum time between flushes \\
      \hline
    \end{tabular}
  \end{center}
\end{table}
At the as-built publish rate (six bursts per node every 5~s, plus one PTP status message per second) the write load is small. The 1~s flush interval keeps PTP status as current as the Node View dashboard, which refreshes every second.

\subsection{QoS Alignment}
Telemetry is published at QoS~0, consistent with the MQTT design. The ingestion pipeline makes no attempt to detect or fill gaps in the historian record resulting from dropped packets or node disconnection. Gaps are visible in the historian and are acceptable for a research testbed.

\subsection{Error Handling}
Telegraf drops malformed payloads that fail JSON parsing and logs an error. No dead-letter queue is implemented. This is acceptable given:
\begin{itemize}[noitemsep]
  \item the two-node testbed operates on a wired LAN with low packet loss probability,
  \item telemetry is QoS~0 by design, so loss is already tolerated at the transport layer, and
  \item a dead-letter mechanism would add operational complexity disproportionate to the research context.
\end{itemize}
